\documentclass[tikz,border=4pt]{standalone}
\usepackage{ctex}
\usepackage{amsmath}
\usetikzlibrary{calc, arrows.meta}

\begin{document}
% part12/01 随机事件：必然/不可能/随机事件示意
\begin{tikzpicture}[scale=1.0, thick]
  % 三类事件分类框
  % 必然事件
  \draw[rounded corners=4pt, fill=green!20, draw=black]
    (-4.5, 0.6) rectangle (-1.5, -0.6);
  \node[font=\small, align=center] at (-3,0.2) {\textbf{必然事件}};
  \node[font=\footnotesize, align=center] at (-3,-0.2) {$P=1$，每次必发生};

  % 不可能事件
  \draw[rounded corners=4pt, fill=pink!30, draw=black]
    (1.5, 0.6) rectangle (4.5, -0.6);
  \node[font=\small, align=center] at (3,0.2) {\textbf{不可能事件}};
  \node[font=\footnotesize, align=center] at (3,-0.2) {$P=0$，每次不发生};

  % 随机事件
  \draw[rounded corners=4pt, fill=orange!20, draw=black]
    (-1.2, 0.6) rectangle (1.2, -0.6);
  \node[font=\small, align=center] at (0,0.2) {\textbf{随机事件}};
  \node[font=\footnotesize, align=center] at (0,-0.2) {$0<P<1$};

  % 概率数轴
  \draw[->] (-5, -1.5) -- (5, -1.5) node[right, font=\small] {$P$};
  \draw (-4.5,-1.4) -- (-4.5,-1.6) node[below, font=\small] {$0$};
  \draw (4.5,-1.4) -- (4.5,-1.6) node[below, font=\small] {$1$};
  \draw (0,-1.4) -- (0,-1.6) node[below, font=\small] {$0.5$};
  % 颜色段
  \fill[pink!30] (-4.5,-1.5) circle(4pt);
  \fill[green!50] (4.5,-1.5) circle(4pt);
  \draw[orange!70, very thick] (-4.5,-1.5) -- (4.5,-1.5);
  \node[above, orange!80!black, font=\footnotesize] at (0,-1.3) {随机事件区间};

  % 示例
  \node[font=\footnotesize, green!60!black, align=left] at (-3,-2.3) {例：太阳从东方升起};
  \node[font=\footnotesize, pink!80!black, align=left] at (3,-2.3) {例：掷骰子得7点};
  \node[font=\footnotesize, orange!80!black, align=center] at (0,-2.3) {例：抛硬币正面朝上};
\end{tikzpicture}
\end{document}
